% !TEX program = pdflatex
% Cahier des charges --- Square One Chess (PFE)
% Compile with pdfLaTeX (twice if you add a TOC later).
% Place this file in docs/ (same folder as logos if you include them).
% =============================================================================
\documentclass[12pt,a4paper,oneside]{article}

\usepackage[T1]{fontenc}
\usepackage[utf8]{inputenc}
\usepackage[french]{babel}
\frenchsetup{AutoSpacePunctuation=false}
\usepackage{lmodern}
\usepackage[a4paper,left=2.5cm,right=2.5cm,top=2.2cm,bottom=2.2cm]{geometry}
\usepackage{scrextend}
\changefontsizes[20pt]{13pt}
\usepackage{graphicx}
\graphicspath{{logos/}{./logos/}{./}}
\usepackage{setspace}
\usepackage[table]{xcolor}
\usepackage{titlesec}
\usepackage{fancyhdr}
\usepackage{float}
\usepackage{booktabs}
\usepackage{array}
\usepackage{tabularx}
\usepackage{longtable}
\usepackage{enumitem}
\usepackage{hyperref}
\usepackage{listings}
\usepackage{caption}

\definecolor{chapterblue}{RGB}{31,72,124}
\definecolor{sectionblue}{RGB}{0,111,192}
\definecolor{subsectionblue}{RGB}{68,113,196}
\definecolor{toclink}{RGB}{0,0,127}
\definecolor{codebg}{RGB}{245,245,245}

\hypersetup{colorlinks=true,linkcolor=toclink,urlcolor=toclink}

\setstretch{1.2}
\setlength{\parindent}{0pt}
\setlength{\parskip}{0.45em}
\setlist[itemize]{leftmargin=1.4em,itemsep=0.25em,topsep=0.3em}
\setlist[enumerate]{leftmargin=1.4em,itemsep=0.25em,topsep=0.3em}

\titleformat{\section}
  {\normalfont\Large\bfseries\color{sectionblue}}
  {\thesection}{0.8em}{}
\titleformat{\subsection}
  {\normalfont\large\bfseries\color{subsectionblue}}
  {\thesubsection}{0.8em}{}
\titleformat{\subsubsection}
  {\normalfont\normalsize\bfseries\color{chapterblue}}
  {\thesubsubsection}{0.8em}{}

\pagestyle{fancy}
\fancyhf{}
\fancyhead[L]{\footnotesize Cahier des charges --- Square One Chess}
\fancyhead[R]{\footnotesize PFE 2025--2026}
\fancyfoot[C]{\thepage}
\renewcommand{\headrulewidth}{0.4pt}
\setlength{\headheight}{14pt}

\newcommand{\StudentName}{Lasmer Dhia Eddin}
\newcommand{\SupervisorName}{M. Aziz Thraya}
\newcommand{\AcademicSupervisor}{M. Hmida Hmida}

\lstset{
  basicstyle=\ttfamily\small,
  backgroundcolor=\color{codebg},
  frame=single,
  breaklines=true,
  columns=fullflexible
}

\begin{document}

% ---------------------------------------------------------------------------
% Page de titre
% ---------------------------------------------------------------------------
\begin{titlepage}
\centering
\vspace*{1cm}

{\Large\bfseries Université de Jendouba}\\[0.3em]
{\large Institut Supérieur d'Informatique du Kef (ISI Kef)}\\[1.2em]

{\color{sectionblue}\rule{\textwidth}{1.2pt}}\\[0.8em]
{\LARGE\bfseries\color{chapterblue} Cahier des charges}\\[0.5em]
{\Large\bfseries Square One Chess}\\[0.3em]
{\large Plateforme web d'échecs en temps réel}\\[0.8em]
{\color{sectionblue}\rule{\textwidth}{1.2pt}}\\[1.2em]

\begin{tabular}{>{\bfseries}l@{\hspace{1em}}p{0.62\textwidth}}
Projet & Projet de fin d'études (PFE) \\[0.35em]
Organisme d'accueil & NCIS (Network and Computer Integration Services), Tunis \\[0.35em]
Étudiant & \StudentName \\[0.35em]
Encadrant industriel & \SupervisorName \\[0.35em]
Encadrant académique & \AcademicSupervisor \\[0.35em]
Version & 1.0 \\
\end{tabular}

\vfill
{\large Année universitaire 2025--2026}
\end{titlepage}

\tableofcontents
\newpage

% ---------------------------------------------------------------------------
\section{Introduction}
% ---------------------------------------------------------------------------

\subsection{Objet du document}
Ce document décrit ce que nous devons réaliser pour le projet \textbf{Square One Chess}.
Il précise le besoin, ce qui est dans le périmètre, les exigences (fonctionnelles et non fonctionnelles), les contraintes techniques et les livrables attendus pour le PFE.

\subsection{Contexte}
Aujourd'hui, beaucoup de joueurs d'échecs utilisent plusieurs outils en même temps:
un site pour jouer, un groupe de discussion pour organiser une partie, un fichier Excel pour noter les résultats, et parfois un autre site pour les puzzles.
Cela complique le suivi du niveau (Elo), de l'historique des parties et de la vie d'un club ou d'une communauté étudiante.

Ce projet a été mené durant un stage à \textbf{NCIS}, une entreprise tunisienne spécialisée dans les services informatiques et le développement d'applications web.

\subsection{Présentation du projet}
\textbf{Square One Chess} est une application web complète (frontend + backend) qui permet:
\begin{itemize}
  \item de créer un compte et de se connecter de manière sécurisée;
  \item de jouer en ligne contre un autre joueur ou contre un moteur d'échecs;
  \item de voir l'historique des parties et de les rejouer;
  \item de s'entraîner avec des puzzles et une revue de partie assistée par l'IA;
  \item de participer à des tournois;
  \item d'ajouter des amis et de discuter via un chat;
  \item de consulter le classement, les statistiques, et (pour un administrateur) les analytics de la plateforme.
\end{itemize}

\subsection{Objectifs}
\begin{table}[H]
\centering
\caption{Objectifs du projet}
\begin{tabularx}{\linewidth}{|>{\bfseries}l|X|}
\hline
\rowcolor{sectionblue}
\textcolor{white}{Type} & \textcolor{white}{\textbf{Objectif}} \\
\hline
Métier & Réunir sur une seule plateforme le jeu, la communauté, les tournois et l'entraînement. \\
\hline
Technique & Mettre en place une architecture claire: interface Angular, API Express et communication temps réel. \\
\hline
Pédagogique & Appliquer Scrum, UML, la sécurité web et le développement TypeScript côté client et serveur. \\
\hline
Qualité & Préparer une démonstration stable et une documentation complète pour la soutenance. \\
\hline
\end{tabularx}
\end{table}

\subsection{Périmètre}

\subsubsection{Ce qui est inclus}
\begin{itemize}
  \item Compte utilisateur: inscription, connexion, session sécurisée (cookie HttpOnly JWT), profil, et selon la configuration: vérification d'e-mail / mot de passe oublié.
  \item Jeu en ligne: lobby, file d'attente, parties en direct (Socket.io), horloge, abandon et nulle.
  \item Partie contre le moteur Stockfish.
  \item Historique, relecture (replay) et revue IA d'une partie terminée.
  \item Puzzles tactiques (quotidien, aléatoire, par thème, plusieurs coups).
  \item Tournois: création, inscription, démarrage, tableau (bracket) et chat.
  \item Amis, présence en ligne et chat.
  \item Classement Elo et statistiques sur le tableau de bord.
  \item Analytics administrateur.
  \item Documentation UML et rapport de projet.
\end{itemize}

\subsubsection{Ce qui n'est pas demandé pour ce PFE}
\begin{itemize}
  \item Application mobile native.
  \item Mode spectateur pour regarder une partie en direct.
  \item Défi entre amis / rematch avancé.
  \item Import de parties PGN externes.
  \item Mise en production à grande échelle (Postgres, Redis, plusieurs serveurs).
  \item Système de modération avancé (signalements, sanctions automatiques).
\end{itemize}

% ---------------------------------------------------------------------------
\section{Étude de l'existant et critique}
% ---------------------------------------------------------------------------

\subsection{Situation actuelle}
Les joueurs amateurs passent souvent d'un outil à l'autre.
Le suivi du classement, la conservation des parties et l'organisation d'une petite communauté restent difficiles sans plateforme dédiée.

\subsection{Limites}
\begin{itemize}
  \item Il n'y a pas une seule application qui couvre le compte, le jeu et l'entraînement.
  \item L'appariement et le contrôle du temps sont peu fiables si tout se fait par messagerie.
  \item Beaucoup de parties sont perdues: difficile de les revoir ensuite.
  \item Le chat et le plateau de jeu ne sont pas liés.
  \item Les tournois sont souvent gérés à la main, avec des risques d'erreurs.
  \item La sécurité des comptes et les rôles (joueur / admin) ne sont pas toujours clairs.
\end{itemize}

\subsection{Solution proposée}
Nous proposons de développer \textbf{Square One}, une plateforme web d'échecs en temps réel.
Elle regroupe l'authentification, le jeu en live, le matchmaking, l'historique, le social, les tournois, les puzzles, la revue IA, les statistiques et un espace analytics pour l'administrateur.

% ---------------------------------------------------------------------------
\section{Acteurs du système}
% ---------------------------------------------------------------------------

\begin{table}[H]
\centering
\caption{Acteurs et interactions}
\begin{tabularx}{\linewidth}{|>{\bfseries}l|X|X|}
\hline
\rowcolor{sectionblue}
\textcolor{white}{Acteur} & \textcolor{white}{\textbf{Description}} & \textcolor{white}{\textbf{Interactions principales}} \\
\hline
Visiteur & Personne qui n'est pas encore connectée & S'inscrire, se connecter \\
\hline
Joueur & Utilisateur connecté & Jouer, s'entraîner, discuter, rejoindre un tournoi, voir ses stats, gérer son profil \\
\hline
Administrateur & Joueur autorisé (liste d'e-mails / identifiants, ou mode démo) & Voir les analytics de la plateforme sur le dashboard \\
\hline
\end{tabularx}
\end{table}

\noindent\textit{Remarque.} Dans l'application, il n'y a pas de page publique pour « parcourir la plateforme » sans compte.
Après connexion, l'utilisateur arrive sur le dashboard.

% ---------------------------------------------------------------------------
\section{Exigences fonctionnelles}
% ---------------------------------------------------------------------------

Les exigences ci-dessous décrivent ce que l'application doit vraiment faire.

\subsection{Compte et sécurité}
\begin{table}[H]
\centering
\caption{Exigences --- compte et sécurité}
\small
\begin{tabularx}{\linewidth}{|c|X|c|}
\hline
\rowcolor{sectionblue}
\textcolor{white}{\textbf{ID}} & \textcolor{white}{\textbf{Exigence}} & \textcolor{white}{\textbf{Priorité}} \\
\hline
EF-01 & Le système permet à un visiteur de s'inscrire (pseudo, e-mail, mot de passe). & Obligatoire \\
\hline
EF-02 & Les mots de passe sont stockés sous forme de hash (bcrypt), jamais en clair. & Obligatoire \\
\hline
EF-03 & Le système permet de se connecter et d'ouvrir une session via cookie HttpOnly JWT. & Obligatoire \\
\hline
EF-04 & Les routes privées exigent une session valide (\texttt{authGuard} / middleware). & Obligatoire \\
\hline
EF-05 & Le joueur peut consulter et modifier son profil (ex.\ pseudo). & Obligatoire \\
\hline
EF-06 & Le joueur peut se déconnecter (invalidation du cookie). & Obligatoire \\
\hline
EF-07 & Vérification d'e-mail et réinitialisation du mot de passe (selon SMTP / variables d'environnement). & Souhaitable \\
\hline
\end{tabularx}
\end{table}

\subsection{Parties en ligne}
\begin{table}[H]
\centering
\caption{Exigences --- parties en ligne}
\small
\begin{tabularx}{\linewidth}{|c|X|c|}
\hline
\rowcolor{sectionblue}
\textcolor{white}{\textbf{ID}} & \textcolor{white}{\textbf{Exigence}} & \textcolor{white}{\textbf{Priorité}} \\
\hline
EF-10 & Le joueur peut rejoindre une file d'attente (contrôle du temps) via Socket.io. & Obligatoire \\
\hline
EF-11 & Le système apparie deux joueurs et crée une partie (\texttt{match:found}). & Obligatoire \\
\hline
EF-12 & Les coups sont validés côté serveur (\texttt{chess.js}) et diffusés en temps réel. & Obligatoire \\
\hline
EF-13 & Une partie gère horloge, abandon, nulle, fin de partie et mise à jour Elo. & Obligatoire \\
\hline
EF-14 & Le joueur peut lancer une partie contre le moteur (difficulté configurable). & Obligatoire \\
\hline
\end{tabularx}
\end{table}

\subsection{Historique et relecture}
\begin{table}[H]
\centering
\caption{Exigences --- historique et relecture}
\small
\begin{tabularx}{\linewidth}{|c|X|c|}
\hline
\rowcolor{sectionblue}
\textcolor{white}{\textbf{ID}} & \textcolor{white}{\textbf{Exigence}} & \textcolor{white}{\textbf{Priorité}} \\
\hline
EF-20 & Le joueur consulte la liste de ses parties terminées. & Obligatoire \\
\hline
EF-21 & Le joueur rejoue une partie coup par coup à partir des coups stockés. & Obligatoire \\
\hline
EF-22 & Le système peut fournir le PGN / détail d'une partie via l'API. & Souhaitable \\
\hline
\end{tabularx}
\end{table}

\subsection{Entraînement}
\begin{table}[H]
\centering
\caption{Exigences --- entraînement}
\small
\begin{tabularx}{\linewidth}{|c|X|c|}
\hline
\rowcolor{sectionblue}
\textcolor{white}{\textbf{ID}} & \textcolor{white}{\textbf{Exigence}} & \textcolor{white}{\textbf{Priorité}} \\
\hline
EF-30 & Le joueur résout des puzzles (quotidien, aléatoire, filtre par thème). & Obligatoire \\
\hline
EF-31 & Résolution multi-coups : coup joueur $\rightarrow$ réponse adverse auto $\rightarrow$ coup suivant ; échec = reset. & Obligatoire \\
\hline
EF-32 & L'API puzzle ne révèle pas la ligne solution au chargement ; l'indice indique la case de départ. & Obligatoire \\
\hline
EF-33 & Les tentatives / réussites sont enregistrées (\texttt{UserPuzzleStats}). & Obligatoire \\
\hline
EF-34 & Le joueur peut ouvrir une revue IA d'une partie terminée (Stockfish / service AI). & Obligatoire \\
\hline
\end{tabularx}
\end{table}

\subsection{Tournois}
\begin{table}[H]
\centering
\caption{Exigences --- tournois}
\small
\begin{tabularx}{\linewidth}{|c|X|c|}
\hline
\rowcolor{sectionblue}
\textcolor{white}{\textbf{ID}} & \textcolor{white}{\textbf{Exigence}} & \textcolor{white}{\textbf{Priorité}} \\
\hline
EF-40 & Un joueur peut créer et rejoindre un tournoi. & Obligatoire \\
\hline
EF-41 & Le créateur peut démarrer le tournoi (si assez de joueurs). & Obligatoire \\
\hline
EF-42 & Le système gère inscriptions, matchs, bracket et résultats liés aux parties. & Obligatoire \\
\hline
EF-43 & Un chat de tournoi est disponible. & Souhaitable \\
\hline
\end{tabularx}
\end{table}

\subsection{Social}
\begin{table}[H]
\centering
\caption{Exigences --- social}
\small
\begin{tabularx}{\linewidth}{|c|X|c|}
\hline
\rowcolor{sectionblue}
\textcolor{white}{\textbf{ID}} & \textcolor{white}{\textbf{Exigence}} & \textcolor{white}{\textbf{Priorité}} \\
\hline
EF-50 & Le joueur gère des demandes d'amis (envoyer, accepter, lister). & Obligatoire \\
\hline
EF-51 & Le joueur communique via chat (global / privé / partie) avec persistance des messages. & Obligatoire \\
\hline
EF-52 & Le système indique la présence en ligne des amis. & Souhaitable \\
\hline
\end{tabularx}
\end{table}

\subsection{Statistiques et administration}
\begin{table}[H]
\centering
\caption{Exigences --- statistiques et administration}
\small
\begin{tabularx}{\linewidth}{|c|X|c|}
\hline
\rowcolor{sectionblue}
\textcolor{white}{\textbf{ID}} & \textcolor{white}{\textbf{Exigence}} & \textcolor{white}{\textbf{Priorité}} \\
\hline
EF-60 & Le joueur consulte le classement (Elo, formats bullet/blitz/rapid selon API). & Obligatoire \\
\hline
EF-61 & Le joueur consulte ses statistiques sur le dashboard. & Obligatoire \\
\hline
EF-62 & L'administrateur consulte des analytics plateforme (utilisateurs, volume de parties, etc.). & Obligatoire \\
\hline
\end{tabularx}
\end{table}

% ---------------------------------------------------------------------------
\section{Exigences non fonctionnelles}
% ---------------------------------------------------------------------------

\begin{table}[H]
\centering
\caption{Exigences non fonctionnelles}
\small
\begin{tabularx}{\linewidth}{|c|>{\bfseries}l|X|}
\hline
\rowcolor{sectionblue}
\textcolor{white}{\textbf{ID}} & \textcolor{white}{\textbf{Catégorie}} & \textcolor{white}{\textbf{Exigence}} \\
\hline
ENF-01 & Sécurité & JWT en cookie HttpOnly ; pas de token en \texttt{localStorage} ; CORS avec credentials ; validation des coups côté serveur. \\
\hline
ENF-02 & Performance & Temps de réponse API acceptable en démo locale ; parties live sans rafraîchir la page. \\
\hline
ENF-03 & Disponibilité & Application utilisable en local pour la soutenance (API + front démarrés). \\
\hline
ENF-04 & Ergonomie & Interface web responsive (Angular + Tailwind), parcours clairs. \\
\hline
ENF-05 & Maintenabilité & Code TypeScript modulaire (routes $\rightarrow$ controllers $\rightarrow$ services $\rightarrow$ Prisma). \\
\hline
ENF-06 & Portabilité & Développement Windows ; stack Node.js ; base SQLite fichier. \\
\hline
ENF-07 & Capacité & Dimensionnement démo : dizaines d'utilisateurs simultanés (SQLite + un process Node) --- limite assumée. \\
\hline
ENF-08 & Traçabilité & Diagrammes UML alignés sur l'implémentation. \\
\hline
\end{tabularx}
\end{table}

% ---------------------------------------------------------------------------
\section{Contraintes}
% ---------------------------------------------------------------------------

\subsection{Techniques}
\begin{table}[H]
\centering
\caption{Contraintes techniques}
\begin{tabularx}{\linewidth}{|>{\bfseries}l|X|}
\hline
\rowcolor{sectionblue}
\textcolor{white}{Élément} & \textcolor{white}{\textbf{Choix}} \\
\hline
Frontend & Angular 17 (SPA), TypeScript, Tailwind, chess.js, Socket.io-client, Chart.js \\
\hline
Backend & Node.js, Express, TypeScript, Socket.io \\
\hline
Persistance & Prisma + SQLite \\
\hline
Moteur & Stockfish (npm) \\
\hline
Auth & JWT + cookie HttpOnly, bcrypt \\
\hline
E-mail & Nodemailer (optionnel selon SMTP) \\
\hline
\end{tabularx}
\end{table}

\subsection{Organisationnelles}
\begin{itemize}
  \item Méthode : \textbf{Scrum} (backlog, sprints, incréments démo).
  \item Durée : stage / PFE selon calendrier ISI Kef et NCIS.
  \item Encadrement : \SupervisorName\ (PO / Scrum Master côté entreprise).
\end{itemize}

\subsection{Juridiques / éthiques}
\begin{itemize}
  \item Pas de stockage de mots de passe en clair.
  \item Respect de la vie privée des comptes de démonstration.
  \item Cohérence entre le code, la documentation et le rapport de stage.
\end{itemize}

\clearpage
% ---------------------------------------------------------------------------
\section{Architecture cible}
% ---------------------------------------------------------------------------

\noindent L'architecture de \textbf{Square One} tient sur une seule vue: le client Angular parle au serveur Node (REST + Socket.io), puis les services utilisent SQLite, Stockfish et éventuellement l'e-mail.

\begin{figure}[H]
\centering
\renewcommand{\arraystretch}{1.35}
\begin{tabular}{|c|}
\hline
\rowcolor{sectionblue!12}
\textbf{Client (navigateur)} \\
Angular 17 + Tailwind + chess.js + Socket.io-client \\
{\footnotesize Port 4200 --- proxy \texttt{/api} et \texttt{/socket.io} en développement} \\
\hline
\multicolumn{1}{c}{$\big\downarrow$ \footnotesize HTTP JSON + cookie JWT \quad|\quad Socket.io (coups, file, chat)} \\
\hline
\rowcolor{sectionblue!12}
\textbf{Serveur Node.js} \\
Express (REST \texttt{/api/*}) \quad+\quad Socket.io (temps réel) \\
{\footnotesize Port 3000 --- Auth JWT dans cookie HttpOnly} \\
\hline
\multicolumn{1}{c}{$\big\downarrow$} \\
\hline
\rowcolor{sectionblue!12}
\textbf{Services métier} \\
Auth \quad Game \quad Puzzle \quad Tournament \quad Stats \quad Social \quad AI \\
\hline
\multicolumn{1}{c}{$\big\downarrow$} \\
\end{tabular}

\vspace{0.4em}
\begin{tabular}{|c|c|c|}
\hline
\rowcolor{sectionblue!8}
\textbf{SQLite} & \textbf{Stockfish} & \textbf{Nodemailer} \\
via Prisma ORM & bot / analyse & e-mail (optionnel) \\
\hline
\end{tabular}
\caption{Architecture de Square One Chess (vue d'ensemble)}
\label{fig:arch-cahier}
\end{figure}

\noindent\textbf{Point important:} le serveur décide si un coup est légal et si la session est valide.
Le frontend affiche le jeu, mais ce n'est pas lui qui fait autorité sur les règles.

\noindent\textbf{Flux principaux (sur la même page):}
\begin{itemize}
  \item \textbf{REST} (\texttt{/api/...}): inscription, connexion, profil, puzzles, tournois, historique, classement.
  \item \textbf{Socket.io}: file d'attente, coups en live, chat et présence.
  \item \textbf{Base SQLite}: utilisateurs, parties, coups, puzzles, tournois, messages.
\end{itemize}

% ---------------------------------------------------------------------------
\section{Cas d'utilisation (synthèse)}
% ---------------------------------------------------------------------------

Le diagramme global des cas d'utilisation (\texttt{docs/uml/01-use-case.puml} / \texttt{fig-usecase-global.png}) résume les interactions suivantes:

\begin{table}[H]
\centering
\caption{Synthèse des cas d'utilisation}
\begin{tabularx}{\linewidth}{|>{\bfseries}l|X|}
\hline
\rowcolor{sectionblue}
\textcolor{white}{Acteur} & \textcolor{white}{\textbf{Cas d'utilisation}} \\
\hline
Visiteur & S'inscrire ; Se connecter \\
\hline
Joueur & Gérer le profil ; Jouer en ligne ; Historique et replay ; Résoudre des puzzles ; Revue d'une partie ; Tournois ; Amis et chat ; Classement ; Statistiques dashboard \\
\hline
Administrateur & Analytics plateforme \\
\hline
\end{tabularx}
\end{table}

% ---------------------------------------------------------------------------
\section{Backlog produit (extrait)}
% ---------------------------------------------------------------------------

Référence complète : \texttt{docs/BACKLOG.md} et \texttt{docs/SPRINT-PLANNING.md}.

\begin{table}[H]
\centering
\caption{Epics et user stories (extrait)}
\begin{tabularx}{\linewidth}{|>{\bfseries}l|X|}
\hline
\rowcolor{sectionblue}
\textcolor{white}{Epic} & \textcolor{white}{\textbf{Exemples d'user stories}} \\
\hline
E1 Compte \& sécurité & US-01 à US-05 \\
\hline
E2 Parties en ligne & US-10 à US-14 \\
\hline
E3 Historique \& replay & US-20, US-21 \\
\hline
E4 Stats \& classement & US-30, US-31 \\
\hline
E5 Social & US-40 à US-42 \\
\hline
E6 Tournois & US-50, US-51 \\
\hline
E7 Puzzles & US-60, US-61 \\
\hline
E8 Coaching IA & US-70 \\
\hline
E9 Administration & US-80 \\
\hline
\end{tabularx}
\end{table}

\noindent Priorisation: P1 = indispensable pour la démo; P2 = important; P3 = amélioration plus tard.

% ---------------------------------------------------------------------------
\section{Critères d'acceptation (recette)}
% ---------------------------------------------------------------------------

Le projet est prêt pour la soutenance si:
\begin{enumerate}
  \item Inscription + connexion fonctionnent et mènent au dashboard (cookie session).
  \item Une partie PvP ou vs bot peut être jouée en live (coups validés, fin de partie).
  \item L'historique et le replay d'une partie terminée fonctionnent.
  \item Au moins un puzzle peut être résolu (y compris ligne multi-coups si proposée).
  \item Classement ou statistiques dashboard accessibles.
  \item Tournoi : création / inscription / démarrage démontrable avec 2 joueurs.
  \item Chat ou amis utilisable depuis l'interface.
  \item Les diagrammes du rapport correspondent aux fonctionnalités livrées.
  \item Le code démarre en local (API \texttt{:3000}, front \texttt{:4200} avec proxy).
\end{enumerate}

% ---------------------------------------------------------------------------
\section{Livrables}
% ---------------------------------------------------------------------------

\begin{table}[H]
\centering
\caption{Livrables du projet}
\begin{tabularx}{\linewidth}{|>{\bfseries}l|X|}
\hline
\rowcolor{sectionblue}
\textcolor{white}{Livrable} & \textcolor{white}{\textbf{Description}} \\
\hline
Code source & Backend + frontend TypeScript (dépôt Git) \\
\hline
Base de données & Schéma Prisma + seed (utilisateurs démo, puzzles) \\
\hline
Documentation technique & README, UML, guides (\texttt{docs/}) \\
\hline
Rapport PFE & Document Latex (\texttt{docs/page-de-garde.tex} et suite) \\
\hline
Cahier des charges & Le présent document \\
\hline
Démo & Scénario de présentation 8--10 min \\
\hline
\end{tabularx}
\end{table}

% ---------------------------------------------------------------------------
\section{Planning indicatif (releases)}
% ---------------------------------------------------------------------------

\begin{table}[H]
\centering
\caption{Releases}
\begin{tabularx}{\linewidth}{|>{\bfseries}c|X|}
\hline
\rowcolor{sectionblue}
\textcolor{white}{Release} & \textcolor{white}{\textbf{Contenu}} \\
\hline
R1 & Compte, auth cookie, profil \\
\hline
R2 & Lobby, matchmaking, parties live, historique, replay \\
\hline
R3 & Social, tournois, puzzles, revue IA, leaderboard, analytics admin \\
\hline
\end{tabularx}
\end{table}

% ---------------------------------------------------------------------------
\section{Risques et hypothèses}
% ---------------------------------------------------------------------------

\begin{table}[H]
\centering
\caption{Risques et mitigations}
\small
\begin{tabularx}{\linewidth}{|X|X|X|}
\hline
\rowcolor{sectionblue}
\textcolor{white}{\textbf{Risque}} & \textcolor{white}{\textbf{Impact}} & \textcolor{white}{\textbf{Mitigation}} \\
\hline
Cookie bloqué (\texttt{localhost} $\neq$ \texttt{127.0.0.1}) & Impossible de rester connecté & Proxy Angular same-origin ; URL \texttt{http://localhost:4200} \\
\hline
SQLite sous charge & Verrous / lenteur & Limiter la démo ; évoquer Postgres en perspectives \\
\hline
Stockfish CPU & Bot / revue lents & Profondeur limitée en démo \\
\hline
SMTP absent & Vérification e-mail & Auto-verify en développement \\
\hline
Puzzles de mauvaise qualité & Démo peu crédible & Seed validé par \texttt{chess.js} \\
\hline
\end{tabularx}
\end{table}

\noindent\textbf{Hypothèses:} on fait la démo en local, avec deux comptes de test prêts, et le jury ouvre l'application dans un navigateur.

% ---------------------------------------------------------------------------
\section{Perspectives (hors cahier initial)}
% ---------------------------------------------------------------------------

\begin{itemize}
  \item Migration PostgreSQL + Redis adapter Socket.io.
  \item Import massif de puzzles (ex.\ dump Lichess filtré).
  \item Mode spectateur, défi ami, notifications.
  \item Application mobile ou PWA.
  \item Tests automatisés (e2e) et CI/CD.
\end{itemize}

% ---------------------------------------------------------------------------
\section{Glossaire}
% ---------------------------------------------------------------------------

\begin{table}[H]
\centering
\caption{Glossaire}
\begin{tabularx}{\linewidth}{|>{\bfseries}l|X|}
\hline
\rowcolor{sectionblue}
\textcolor{white}{Terme} & \textcolor{white}{\textbf{Définition}} \\
\hline
FEN & Notation d'une position d'échiquier \\
\hline
SAN & Notation algébrique d'un coup (\texttt{Nf3}, \texttt{O-O}\ldots) \\
\hline
JWT & Jeton signé d'identité \\
\hline
HttpOnly cookie & Cookie non accessible en JavaScript \\
\hline
Elo & Système de classement \\
\hline
Matchmaking & Appariement automatique de joueurs \\
\hline
Socket.io & Communication temps réel bidirectionnelle \\
\hline
Prisma & ORM TypeScript vers la base \\
\hline
\end{tabularx}
\end{table}

% ---------------------------------------------------------------------------
\section{Approbation}
% ---------------------------------------------------------------------------

\begin{table}[H]
\centering
\caption{Fiche d'approbation}
\begin{tabularx}{\linewidth}{|X|X|c|c|}
\hline
\rowcolor{sectionblue}
\textcolor{white}{\textbf{Rôle}} & \textcolor{white}{\textbf{Nom}} & \textcolor{white}{\textbf{Date}} & \textcolor{white}{\textbf{Signature}} \\
\hline
Étudiant & \StudentName & & \\[2.2em]
\hline
Encadrant industriel & \SupervisorName & & \\[2.2em]
\hline
Encadrant académique & \AcademicSupervisor & & \\[2.2em]
\hline
\end{tabularx}
\end{table}

\vfill
\begin{center}
\textit{Ce cahier des charges correspond à l'application Square One Chess livrée pour le PFE.}
\end{center}

\end{document}
